\subsection{Evolution of reduced-order solutions without conserved quantities}

Consider the map:
\begin{equation}
    \begin{aligned}
        u: D\times \mathbb{R}^+ & \rightarrow \mathbb{R}^p \\
        (\mathbf{x}, t) & \mapsto u(\mathbf{x}, t)
    \end{aligned}
\end{equation}
where $D \subset \mathbb{R}^k$ is the spatial domain. 
For simplicity, the following work is restricted to the case $p=1$.
This map is governed by a PDE of the form:
\begin{equation}
    \label{eq:pde}
    \begin{aligned}
        & \pd{u}{t} = F(u) & \text{ in } D\times \mathbb{R}^+  , \\
        & u(\mathbf{x}, 0) = u_0(\mathbf{x}),  & \text{ in } D,
    \end{aligned}
\end{equation}
where $F$ is a potentially nonlinear differential operator and $u_0$ is the initial data.

An \textit{ansatz} $\hat{u}(\mathbf{x}, \mathbf{q}(t))$ is constructed to be the approximate solution to \eqref{eq:pde}. 
The functional dependence on the spatial variable $x$ is prescribed a priori, levaraging proprieties of the PDE. 

% An approximate solution of the form $\hat{u}(\mathbf{x}, \mathbf{q}(t))$ to \eqref{eq:pde} is sought, which will be refered as \textit{ansatz}.
% The ansatz $\hat{u}(\mathbf{x}, \mathbf{q}(t))$ depe

The ansatz $\hat{u}(\mathbf{x}, \mathbf{q}(t))$ is seen as a map from the parameters $\mathbf{q}\in \Omega \subset \mathbb{R}^n$ to the function space H:
\begin{equation}
    \label{eq:ansatz}
    \begin{aligned}
        \hat{u} : \Omega & \rightarrow H \\
        \mathbf{q} & \mapsto \hat{u}(\cdot, \mathbf{q})
    \end{aligned}
\end{equation}

\begin{defs}[Ansatz Manifold]
    The Ansatz Manifold is defined as 
    \begin{equation}
        \mathcal{M} = \{ u \in H : \exists \mathbf{q} \in \Omega \text{ with } u = \hat{u}(\cdot , \mathbf{q})  \} \subset H
    \end{equation}
    which is the image of \eqref{eq:ansatz}.
\end{defs}

\begin{assum}
    \label{assum:1}
    Assume that the map $\hat{u}: \Omega \rightarrow H$ has the following properties:
    \begin{enumerate}
        \item it is injective and at least once continuously with respect to $q$;
        \item for every $\mathbf{q}\in \Omega$, the ansatz is $l$-differentiable with respect to $\mathbf{x}$ where $l \in \mathbb{N}$ is the highest oder of spatial derivatives appearing in \eqref{eq:pde};
        \item the map $\hat{u}: \Omega \rightarrow H$ is an immersion.
    \end{enumerate}
\end{assum}

\begin{equation}
    T_{\hat{u}}\mathcal{M} := \text{span}\left\{\pd{\hat{u}}{q_1}, \pd{\hat{u}}{q_2}, ..., \pd{\hat{u}}{q_n}\right\}
\end{equation}